% Compiled by `doit compile_report` (latexmk -xelatex). Every number and exhibit
% below is generated by code into ../_output and pulled in with \input or
% \includegraphics. Never type a statistic by hand.
\documentclass[11pt]{article}
\usepackage[margin=1in]{geometry}
\usepackage{booktabs}
\usepackage{graphicx}
\usepackage{hyperref}

% These paths assume the default OUTPUT_DIR (_output). If you relocate it in
% .env, update \PathToOutput here and the paths in chartbook.toml too.
\newcommand*{\PathToAssets}{../assets}
\newcommand*{\PathToOutput}{../_output}
\input{\PathToOutput/stats.tex}

\title{The U.S. Treasury Term Spread: A Reproducible Pipeline}
\author{Your Name}
\date{\today}

\begin{document}
\maketitle

\section{Overview}
This report is rebuilt from raw data with a single command (\texttt{doit}).
It uses the zero-coupon Treasury yield curve of G\"urkaynak, Sack and Wright
(2007), published daily by the Federal Reserve. The cleaned sample runs from
\SampleStart{} to \SampleEnd{} and covers \NumDates{} trading days. Over this
period the 10-year minus 2-year spread averaged \MeanSpread{} percentage points,
and the curve was inverted on \ShareInverted\% of days.

\section{Data and Validation}
Table~\ref{tab:waterfall} lists every cleaning filter in the order applied and
what it removed. Automated tests confirm that the panel has one row per date
and maturity, that the 2-year and 10-year yields exist on every date, and that
no week passes without data.

\begin{table}[htbp]
  \centering
  % A caption states what the reader should learn, not just what the table is.
  % Rewrite the bracketed text once you have looked at the generated numbers.
  \caption{Cleaning waterfall. [Takeaway, e.g.: the sample-window filter removes
           the pre-\SampleStart{} history; the implausible-value filter removes
           almost nothing, so unit errors are not a concern.]}
  \label{tab:waterfall}
  \input{\PathToOutput/waterfall_table.tex}
\end{table}

\section{Results}
\begin{table}[htbp]
  \centering
  \caption{Zero-coupon yields by maturity (percent). [Takeaway, e.g.: how the
           average and the volatility of yields change with maturity.]}
  \label{tab:summary}
  \input{\PathToOutput/summary_table.tex}
\end{table}

\begin{figure}[htbp]
  \centering
  \includegraphics[width=0.85\linewidth]{\PathToOutput/term_spread.png}
  \caption{The 10-year minus 2-year spread. [Takeaway, e.g.: when inversions
           (below zero) occur and how long they last.]}
  \label{fig:spread}
\end{figure}

\end{document}
